\begin{abstract}
\textbf{Objective:} We investigate how a
mischaracterized covariance validation
objective changes the regularisation
selected for cross-subject EEG decoding,
and whether alternative source-only
selection criteria improve held-person
discrimination or downside risk.
\textbf{Methods:} In the original
source-grouped covariance/tangent-LDA
pipeline, we distinguish validation
class-centred Gaussian covariance
scores, source-trained-mean predictive
Gaussian scores and source-person
AUC. We derive an exact class-mean-shift
penalty and source-only regularisation
selection breakpoints. We preserve
original prospective EEG results,
evaluate corrected estimands in two
previously exposed public cohorts,
and conduct repository-protocol-locked
new runs on independent 22-channel and
3-channel EEG datasets. The latter
additionally tests a fixed worst-quarter
source-person AUC objective (CVaR).
We also execute an exact positive-negative
pair rank-transition decomposition on
held source-validation people; this
mechanistic diagnostic is not a new
target-risk predictor.
\textbf{Results:} The original EEGMMIDB
block ($N=20$) gave $+0.0561$ mean AUC
for discrimination over the historical
class-centred criterion
(Holm $p=0.00597$), with five harmed
individuals; two originally frozen
BNCI tests were not individually
significant and a prospective
filter-bank hypothesis failed.
Retrospective proper Gaussian
comparators nearly tied source
discrimination on Shin2017A and
trailed by 0.00682 AUC on
BNCI2014\_002. A newly protocol-locked
22-channel BNCI2014\_001 cohort
($N=9$) yielded 0.78905 and 0.78341
AUC for source discrimination and
correct Gaussian prediction.
A separate low-rank BNCI2014\_004
cohort ($N=9$, 6,520 trials,
3 channels) yielded AUCs of
0.77800, 0.77658 and 0.77714
for discrimination, proper predictive
selection and source-CVaR, respectively.
The risk-oriented comparator changed
shrinkage decisions but did not improve
mean AUC or the adverse-subject count.
On BNCI2014\_002, the 14 retrospective
source-pool audits showed mean source
AUC difference +0.02680 for source
macro versus historical class-centred
selection with 15.31\% gross trial-pair
rank-state churn. These are SOURCE-only
measurements, not new target improvements.
\textbf{Conclusion:} The source
validation estimand changes EEG
regularisation decisions across
acquisition geometries, but there
is no confirmed universally superior
selector. Reporting the actual
objective, adverse transfer and
source-target information boundary
is essential before generalized
EEG decoding claims.
\end{abstract}
